% SPDX-License-Identifier: CC-BY-SA-4.0
% Author: Matthieu Perrin
% Part: <Nom de la partie>
% Section: <Nom de la section>
% Sub-section: <Nom de la sous-section>  % (facultatif, laisser vide si non utilisé)
% Frame: <Titre de la slide>

\begingroup

\begin{frame}{Un ordinateur est une machine programmable}

  \begin{block}{Quatre machines algorithmiques}
    Ces quatre machines exécutent un \structure{algorithme} modélisable par une MTD
    \begin{description}
    \item[Machine dédiée :] réalise une famille fixe de calculs (Anticythère, Pascaline)
    \item[Machine programmable :] le comportement dépend d'un programme (Zuse 1)
    \item[Machine universelle :] peut simuler n'importe quelle autre machine
    \end{description}
  \end{block}
      
  \begin{block}{Un ordinateur est une machine universelle}
    L'algorithme implémenté par un ordinateur résout le problème suivant :
    \Probleme{Universal}{
      \begin{itemize}
      \item un \structure{programme} qui décrit un algorithme $A$
      \item une \structure{instance} $u$ \alert{de $A$}
      \end{itemize}
    }{
      Quel résultat retourne l'exécution de $A$ sur $u$ ? 
    }
    \begin{itemize}
      \item Un ordinateur peut \structure{simuler} n'importe quelle machine.
      \item Un \structure{programme} est une \alert{donnée} comme les autres
    \end{itemize}
  \end{block}
  
\end{frame}

\endgroup
